\documentclass[10pt,a4paper]{amsart}
\usepackage{amsmath,amssymb,amsthm,graphicx}
\usepackage{verbatim}
\newenvironment{reference}{\comment}{\endcomment}
%\newenvironment{reference}{}{}
\newenvironment{slogan}{\comment}{\endcomment}
\newenvironment{history}{\comment}{\endcomment}

% For commutative diagrams we use Xy-pic
\usepackage[all]{xy}

% We use 2cell for 2-commutative diagrams.
\xyoption{2cell}
% Theorem environments.
%
\theoremstyle{plain}
\newtheorem{theorem}[subsection]{Theorem}
\newtheorem{proposition}[subsection]{Proposition}
\newtheorem{lemma}[subsection]{Lemma}

\theoremstyle{definition}
\newtheorem{definition}[subsection]{Definition}
\newtheorem{example}[subsection]{Example}
\newtheorem{exercise}[subsection]{Exercise}
\newtheorem{situation}[subsection]{Situation}

\theoremstyle{remark}
\newtheorem{remark}[subsection]{Remark}
\newtheorem{remarks}[subsection]{Remarks}

\numberwithin{equation}{subsection}

% Macros
%
\def\lim{\mathop{\mathrm{lim}}\nolimits}
\def\colim{\mathop{\mathrm{colim}}\nolimits}
\def\Spec{\mathop{\mathrm{Spec}}}
\def\Hom{\mathop{\mathrm{Hom}}\nolimits}
\def\Ext{\mathop{\mathrm{Ext}}\nolimits}
\def\SheafHom{\mathop{\mathcal{H}\!\mathit{om}}\nolimits}
\def\SheafExt{\mathop{\mathcal{E}\!\mathit{xt}}\nolimits}
\def\Sch{\mathit{Sch}}
\def\Mor{\mathop{\mathrm{Mor}}\nolimits}
\def\Ob{\mathop{\mathrm{Ob}}\nolimits}
\def\Sh{\mathop{\mathit{Sh}}\nolimits}
\def\NL{\mathop{N\!L}\nolimits}
\def\CH{\mathop{\mathrm{CH}}\nolimits}
\def\proetale{{pro\text{-}\acute{e}tale}}
\def\etale{{\acute{e}tale}}
\def\QCoh{\mathit{QCoh}}
\def\Ker{\mathop{\mathrm{Ker}}}
\def\Im{\mathop{\mathrm{Im}}}
\def\Coker{\mathop{\mathrm{Coker}}}
\def\Coim{\mathop{\mathrm{Coim}}}

% Boxtimes
%
\DeclareMathSymbol{\boxtimes}{\mathbin}{AMSa}{"02}

%
% Macros for moduli stacks/spaces
%
\def\QCohstack{\mathcal{QC}\!\mathit{oh}}
\def\Cohstack{\mathcal{C}\!\mathit{oh}}
\def\Spacesstack{\mathcal{S}\!\mathit{paces}}
\def\Quotfunctor{\mathrm{Quot}}
\def\Hilbfunctor{\mathrm{Hilb}}
\def\Curvesstack{\mathcal{C}\!\mathit{urves}}
\def\Polarizedstack{\mathcal{P}\!\mathit{olarized}}
\def\Complexesstack{\mathcal{C}\!\mathit{omplexes}}
% \Pic is the operator that assigns to X its picard group, usage \Pic(X)
% \Picardstack_{X/B} denotes the Picard stack of X over B
% \Picardfunctor_{X/B} denotes the Picard functor of X over B
\def\Pic{\mathop{\mathrm{Pic}}\nolimits}
\def\Picardstack{\mathcal{P}\!\mathit{ic}}
\def\Picardfunctor{\mathrm{Pic}}
\def\Deformationcategory{\mathcal{D}\!\mathit{ef}}
\usepackage{cleveref}
\crefname{originallemma}{original lemma}{original lemmas}
\Crefname{originallemma}{Original lemma}{Original lemmas}
\makeatletter
\newlabel{flat-lemma-flat-dimension-n}{{05MV}{source}{}{sourcepointer-1}{}}
\newlabel{flat-lemma-flat-dimension-n@cref}{{[originallemma][0][]05MV}{[1][0][]source}{}{}{}}
\makeatother
\crefname{originalreference}{original reference}{original references}
\Crefname{originalreference}{Original reference}{Original references}
\makeatletter
\newlabel{lemma-pre-flat-dimension-n}{{0CWF}{source}{}{sourcepointer-2}{}}
\newlabel{lemma-pre-flat-dimension-n@cref}{{[originalreference][0][]0CWF}{[1][0][]source}{}{}{}}
\makeatother
\crefname{originalreference}{original reference}{original references}
\Crefname{originalreference}{Original reference}{Original references}
\makeatletter
\newlabel{situation-flat-dimension-n}{{05MT}{source}{}{sourcepointer-3}{}}
\newlabel{situation-flat-dimension-n@cref}{{[originalreference][0][]05MT}{[1][0][]source}{}{}{}}
\makeatother
\title{The functor of flatness in dimensions at least n is an fpqc sheaf and is limit preserving}
\author{The Stacks Project Authors}
\date{stacks-project:spaces-flat.tex:local-snapshot-20261012; GNU FDL 1.2 or later}
\begin{document}
\maketitle
\noindent\textit{Editorial note.} All delivered mathematical text is the author's own, copied from the pinned Stacks Project chapter file spaces-flat.tex; no statement or proof step was rewritten, reordered or removed.\par\medskip
\noindent\textit{Editorial note.} The counted proof is the author's complete argument: a chart reduction identifying the functor with a quasi-affine chart functor, applied to base changed fpqc coverings for the sheaf property and to a manufactured filtered limit for limit preservation\par\medskip
\noindent\textit{Editorial note.} The label lemma-flat-dimension-n names the counted statement, and the proof environment immediately following it in the source is delivered verbatim as its proof.\par\medskip
\noindent\textit{Editorial note.} The counted statement is carried by the author's own lemma environment in the chapter source.\par\medskip
\noindent\textit{Allowed standard prerequisites (editorial).}\begin{itemize}
\item The situation defining the functor classifying the points where a module is flat in dimensions at least n, and the chart lemma reducing the functor to a quasi-affine situation over the base
\item The sheaf property for the fpqc topology and limit preservation in that quasi-affine situation
\item The behaviour of fpqc coverings and of filtered limits of affine schemes under base change
\end{itemize}\medskip
\ifdefined\hypertarget\hypertarget{sourcepointer-1}{}\fi
\noindent\textit{Source reference.} Original lemma 05MV: sources/stacks-chow/tags.txt, source lines 7284--7284. Prerequisite pointer invoked at proof line(s) 2139; it names a standard fact of a sibling chapter that the counted argument uses, and every step of the counted proof built on it is retained verbatim in the delivered body.\par\medskip
\ifdefined\hypertarget\hypertarget{sourcepointer-2}{}\fi
\noindent\textit{Source reference.} Original reference 0CWF: sources/stacks-spaces-flat/pinned-original-source.tex, source lines 2030--2030. Prerequisite pointer invoked at proof line(s) 2136; it names a standard fact of the same chapter that the counted argument uses, and every step of the counted proof built on it is retained verbatim in the delivered body.\par\medskip
\ifdefined\hypertarget\hypertarget{sourcepointer-3}{}\fi
\noindent\textit{Source reference.} Original reference 05MT: sources/stacks-spaces-flat/pinned-original-source.tex, source lines 2085--2085. Prerequisite pointer invoked at proof line(s) 2123; it names a standard fact of the same chapter that the counted argument uses, and every step of the counted proof built on it is retained verbatim in the delivered body.\par\medskip
\begin{lemma}
\label{lemma-flat-dimension-n}
In Situation \ref{situation-flat-dimension-n}.
\begin{enumerate}
\item The functor $F_n$ satisfies the sheaf property for the fpqc topology.
\item If $f$ is quasi-compact and locally of finite presentation
and $\mathcal{F}$ is of finite presentation, then the functor $F_n$ is
limit preserving.
\end{enumerate}
\end{lemma}

\begin{proof}
Proof of (1). Suppose that $\{T_i \to T\}$ is an fpqc covering of
a scheme $T$ over $Y$. We have to show that if $F_n(T_i)$ is nonempty
for all $i$, then $F_n(T)$ is nonempty.
Choose a diagram as in part (1) of Lemma \ref{lemma-pre-flat-dimension-n}.
Denote $F'_n$ the corresponding functor for
$\varphi^*\mathcal{F}$ and the morphism $U \to V$.
By More on Flatness, Lemma \ref{flat-lemma-flat-dimension-n}
we have the sheaf property for $F'_n$.
Thus we get the sheaf property for $F_n$ because
for $T \to Y$ we have $F_n(T) = F'_n(V \times_Y T)$
by Lemma \ref{lemma-pre-flat-dimension-n}
and because $\{V \times_Y T_i \to V \times_Y T\}$
is an fpqc covering.

\medskip\noindent
Proof of (2). Suppose that $T = \lim_{i \in I} T_i$ is a filtered limit
of affine schemes $T_i$ over $Y$ and assume that $F_n(T)$ is nonempty.
We have to show that $F_n(T_i)$ is nonempty for some $i$.
Choose a diagram as in
part (1) of Lemma \ref{lemma-pre-flat-dimension-n}.
Fix $i \in I$ and choose an affine open $W_i \subset V \times_Y T_i$
mapping surjectively onto $T_i$. For $i' \geq i$ let $W_{i'}$
be the inverse image of $W_i$ in $V \times_Y T_{i'}$ and
let $W \subset V \times_Y T$ be the inverse image of $W_i$.
Then $W = \lim_{i' \geq i} W_i$ is a filtered limit of affine
schemes over $V$. By
Lemma \ref{lemma-pre-flat-dimension-n} again
it suffices to show that $F'_n(W_{i'})$ is nonempty for
some $i' \geq i$. But we know that $F'_n(W)$ is nonempty
because of our assumption that $F_n(T) = F'_n(V \times_Y T)$
is nonempty. Thus we can apply
More on Flatness, Lemma \ref{flat-lemma-flat-dimension-n}
to conclude.
\end{proof}

\end{document}
